% Generated draft: requires mathematical and source review.
\documentclass[11pt]{article}
\usepackage[margin=25mm]{geometry}
\usepackage{fontspec}
\setmainfont{Latin Modern Roman}
\usepackage{amsmath,amssymb,mathtools,unicode-math}
\setmathfont{Latin Modern Math}
\usepackage{enumitem,needspace,xurl,hyperref}
\setlength{\parindent}{0pt}
\setlength{\parskip}{5pt}
\setlength{\emergencystretch}{3em}
\newcommand{\N}{\mathbb{N}}
\newcommand{\Z}{\mathbb{Z}}
\newcommand{\Q}{\mathbb{Q}}
\newcommand{\R}{\mathbb{R}}
\newcommand{\C}{\mathbb{C}}
\newcommand{\F}{\mathbb{F}}
\newcommand{\A}{\mathbb{A}}
\newcommand{\PP}{\mathbb{P}}
\newcommand{\E}{\mathbb{E}}
\newcommand{\eps}{\varepsilon}
\newcommand{\dd}{\,\mathrm{d}}
\begin{document}
\section*{UnsolvedMath ID: 30004821}
\textbf{Source problem code:} catalogue entry
\subsection{Definitions and mathematical statement}
Write \(\N_{>0}=\{1,2,\ldots\}\) and let
\[
\operatorname{Inc}(\N_{>0})=\{\pi:\N_{>0}\longrightarrow\N_{>0}:\pi(i)<\pi(j)\text{ whenever }i<j\}
\]
be the monoid of strictly increasing maps, with composition as its operation. For an \(R_n\)-module \(M\), its projective dimension \(\operatorname{pd}_{R_n}(M)\) is the least length of a projective resolution, or \(\infty\) if no finite projective resolution exists.

The supplied record does not specify the rings \(R_n\), their coefficient field, their variables or grading, their transition maps, or the action of \(\operatorname{Inc}(\N_{>0})\). Consequently, the conjecture below is relative to the intended ambient system, which must be recovered from the cited source before a complete formalization is possible. Under the usual compatible-map interpretation, an increasing map \(\pi\) with \(\pi(m)\leq n\) induces a ring homomorphism \(\pi_{m,n}:R_m\to R_n\), and invariance means
\[
\pi_{m,n}(I_m)\subseteq I_n
\qquad(m,n\geq1,\ \pi\in\operatorname{Inc}(\N_{>0}),\ \pi(m)\leq n).
\]
This is an explicitly conditional interpretation; the record does not supply the action needed to instantiate it or the precise convention for ``filtration.''

\textbf{Conjecture 3, attributed in the supplied excerpt to [3, Conjecture 1.3].} For the intended ring system and action, let \(I=(I_n)_{n\geq1}\) be any \(\operatorname{Inc}(\N_{>0})\)-invariant filtration such that \(I_n\) is a homogeneous ideal and
\[
0\neq I_n\subsetneq R_n\qquad\text{for every }n\geq1.
\]
Is the sequence \(\bigl(\operatorname{pd}_{R_n}(R_n/I_n)\bigr)_{n\geq1}\) eventually linear? Explicitly, the target is
\[
\exists a\in\Z_{\geq0}\ \exists b\in\Z\ \exists N\in\N_{>0}\ \forall n\geq N:\qquad
\operatorname{pd}_{R_n}(R_n/I_n)=an+b.
\]
Here eventual linearity means eventual equality with an affine function, with finite values; integrality and nonnegativity of the slope follow from the nonnegative integer values at consecutive sufficiently large indices.

The supplied excerpt reports that Conjecture 3 is proved for saturated chains of monomial ideals. Neither ``saturated'' nor the precise ambient monomial setting is defined in this record, so this reported special case cannot be formulated completely here. The excerpt separately reports an asymptotic result for Conjecture 2 for all chains of monomial ideals, saturated or otherwise; Conjecture 2 is not supplied, and that report must not be substituted for the present conjecture. These are supplied source claims, not independently verified results.
\subsection{Short English statement}
For every invariant filtration of nonzero proper homogeneous ideals in the intended ring system, the conjecture asks whether the projective dimensions of the quotients eventually equal \(an+b\). The ambient rings and action are missing from the record. Supplied dated curation labels the problem partially solved and reports the saturated monomial case as proved; human review must verify the setting and these claims.
\subsection{Sources}
\begin{itemize}
\item Local UnsolvedMath registry: supplied record with permanent ID \(30004821\), title ``Eventual Linearity of Equivariant Projective Dimensions,'' exact source code \texttt{catalogue entry}, in the Oberwolfach Reports Open Problems collection.
\item Supplied source citation: \emph{MFO-RIMS Tandem Workshop: Symmetries on Polynomial Ideals and Varieties (2022)}, Oberwolfach Reports, DOI 10.4171/owr/2021/43, pp.\ 2365--2396. \url{https://doi.org/10.4171/owr/2021/43}. The original excerpt labels the assertion Conjecture 3 and attributes it to ``[3, Conjecture 1.3]''; the bibliographic entry numbered [3] is not supplied. No individual proposer is identified.
\item Supplied literature source labeled ``Le--Nguyen'': \url{https://arxiv.org/abs/2206.15141}. The record does not provide its title or a precise theorem reference.
\item Literature assessment dated exactly \texttt{2026-08-22}: supplied dated curation, not independently verified. It labels the problem partially solved, reports a proof for saturated chains of monomial ideals, and states that the claim for arbitrary nonzero proper homogeneous invariant filtrations is not settled. Its listed sources are the OWR report and Le--Nguyen links above. The background requests independent verification of time-sensitive status and cited claims.
\item The registry supplies the statement-status label \texttt{corrected\_verified} and says that the official report was checked directly. This is a supplied verification claim; no independent source consultation is asserted here. The proposed year and year tag are 2021, whereas the supplied workshop citation includes 2022; both are retained without resolving their relationship.
\end{itemize}
\end{document}
